\documentclass[12pt, ukrainian]{article}
\usepackage[a4paper]{geometry}
\setlength\parindent{0mm}
\geometry{top=1.1cm,bottom=1.1cm,left=1.1cm,right=1.1cm}
\pagestyle{empty}
\usepackage[T2A]{fontenc}
\usepackage[utf8]{inputenc}
\usepackage[ukrainian]{babel}
\usepackage{graphicx}
\usepackage{caption}
\DeclareCaptionFormat{nocaption}{}
\captionsetup{format=nocaption,aboveskip=0pt,belowskip=0pt}
\usepackage[Export]{adjustbox}
\adjustboxset{max size={0.9\linewidth}{0.9\paperheight}}
\begin{document}
{{BETA 21.11.2024 Контрольна робота \No 2, варіант  \No \textbf { 60 }\par}}
{
%\begin {large}

{%Завдання 1-10: Що буде виведено за виконання (повідомлення інтерпретатора про помилку %позначати error)?

%Завдання 11-12 (на звороті): написати функцію для розв'язання задачі. Оцінюється економність обчислень та якість коду.

%Тематика задач: використання циклів, програмування з використанням рекурентних співвідношень,
%обробка списків та рядків.

\vspace{5mm} \begin{minipage}[t]{8.5cm}
